\section{Introduction}
\label{sec:intro}

Procurement authorities need triage tools that can flag suspicious
environments with minimal data. Most existing bid-rigging screens
require granular bid amounts---within-tender variance, kurtosis,
and coefficient of variation
\citep{imhof2019detecting, huber2019machine}---yet many procurement
systems record only who bid and who won. When cover bidders
coordinate their bids tightly around the designated winner,
dispersion \emph{falls} and these screens lose power. The
enforcement gap is practical: how should an agency with
participation records but no bid microdata decide where to direct
its investigative resources?

We propose a screen built on a simpler footprint. Cover bidding
requires firms to show up and lose---repeatedly. A firm that bids
on 50 tenders over a decade without winning a single one is either
extraordinarily unlucky or engaged in behavior that warrants
scrutiny. We call these firms \emph{frequent losers} (FL) and show
that this participation anomaly---observable from win/loss records
alone---flags procurement environments with systematically higher
conditional prices and strong proximity to known cartel
environments. A simple conceptual framework
(Section~\ref{sec:structural}) organizes the screening intuition by
distinguishing two cover-bidding regimes, but the screen's value
rests on its empirical performance, not on the framework.

We test this logic on S\~ao Paulo's Bolsa Eletr\^onica de Compras
(BEC), covering 4.5 million tender-items and 40 million bids over
2009--2019. Among 41,000 firms, 16,843 never win a single tender;
a simple IQR-based outlier rule identifies 2,735 whose
participation frequency is difficult to reconcile with
profit-maximizing entry. The FL rule is a first-stage triage device
that prioritizes suspicious environments for investigation, not a
tool for establishing cartel membership or adjudicating liability.

Against CADE (Brazil's competition authority), the screen achieves
AUC~$= 0.94$---higher than a CV-based proxy for
\citet{imhof2019detecting}-style screens (AUC~$= 0.79$), though the
comparison uses a coarser implementation than the full Imhof pipeline
(Section~\ref{sec:results_detection}). FL firms co-participate with
convicted cartelists at 3.5$\times$ the baseline rate
($p < 0.001$). Three CADE-convicted cover bidders are themselves
classified as FL (Section~\ref{sec:cade}).

FL-flagged tenders show 3.6--7.7\% higher conditional prices in
within-item, within-year, within-purchasing-unit comparisons. The
range spans temporal cross-fitting (3.6\%), OLS with high-dimensional
fixed effects (6.4\%), and matching (CEM: 7.7\%; IPW: 5.5\%). After
absorbing item fixed effects the raw gap collapses from 2.43 to 0.050
log points (overlap 0.978; Table~\ref{tab:conditional_descstats},
Appendix~\ref{app:competition}), so the regressions compare
observationally similar tenders within tight product categories.
Sensitivity analysis (\citealt{cinelli2020making},
$RV = 17.5\%$) sets a high bar for omitted confounders.

Five supporting diagnostics assess whether the screen's behavior is
coherent with coordinated cover bidding: the conditional price
association concentrates in competitive markets, Bajari--Ye tests
reject exchangeability of FL and non-FL bid residuals, FL--winner
pairs recur beyond chance, the FL coefficient is larger where cover
bidding is voluntary, and a horse-race regression shows FL captures
information nearly orthogonal to bid-level screens (correlation
0.06). No single diagnostic is dispositive; the joint pattern
strengthens the rationale for deploying the screen.

\subsection{Contributions}

\emph{First}, we build a participation-based screen that requires
only win/loss records---no bid microdata, no enforcement priors, no
supervised training labels---and is immediately deployable wherever
procurement systems record who bid and who won.

\emph{Second}, the screen achieves AUC~$= 0.94$ against
competition-authority convictions, flags environments with
3.6--7.7\% higher conditional prices, and complements bid-level
tools that capture different collusion signatures (correlation 0.06).

\emph{Third}, five independent diagnostics---Bajari--Ye
bid-coordination tests, network-split heterogeneity, regime
selection, dyadic linkage, and invitation variation---are jointly
coherent with coordinated cover bidding, providing supporting
evidence for the screen's economic rationale.

\medskip
\noindent\textbf{Non-claims.}\quad Three boundaries define what this
paper does \emph{not} claim. The conditional price association is
not a causal estimate of the effect of collusion on prices. The FL
classification identifies a participation anomaly, not cartel
membership: it flags suspicious environments, not guilty firms.
The diagnostics are coherent with coordinated cover bidding but do
not uniquely identify the mechanism and do not constitute
sanction-grade evidence. The paper's contribution is the screen
itself---its construction, external validation, and the coherent
diagnostic pattern that supports its deployment as an investigative
prioritization tool.

\medskip
\noindent\textbf{Scope.}\quad The remainder of the paper proceeds
as follows. Section~\ref{sec:literature} reviews related work.
Section~\ref{sec:data} describes the data and FL definition.
Section~\ref{sec:structural} presents the conceptual framework.
Section~\ref{sec:empirical} lays out the empirical strategy.
Section~\ref{sec:cade} reports the CADE consistency check.
Sections~\ref{sec:results}--\ref{sec:mechanisms} present results and
diagnostics. Section~\ref{sec:robustness} reports robustness checks.
Section~\ref{sec:limitations} discusses limitations.
Section~\ref{sec:conclusion} proposes a three-stage enforcement
pathway and discusses portability.
